Natural-language interfaces for decentralized finance (DeFi) promise to
make multi-step protocol interactions accessible to non-expert users, but
they are typically evaluated only on whether the generated workflow graph
looks structurally valid. We argue that this is the wrong bar for a
high-stakes execution domain, and we study a stricter question: are
generated DeFi workflows \emph{safely executable}? We introduce
\textsc{DeFiFlowBench}, a benchmark of 120 expert-authored prompts spanning
swaps, limit orders, cross-chain transfers, and compositional tasks, with
difficulty tiers, paraphrase clusters, and adversarial cases. Each generated
workflow is scored on three nested static layers---structural validity,
executable configuration completeness, and declared safety---and on an
execution layer that runs each generated swap on a real EVM. To make scores
interpretable we bound the scale with an oracle solver (which scores $1.00$
on every layer and is on-chain safe) and null and random-node baselines
(which score $0.00$), and we validate the static safety proxy against on-chain
outcomes: on our data the proxy never certifies a workflow that then executes
unsafely. Evaluating a deterministic template, a deployed no-code DeFi
generator, and direct, constrained, few-shot, and safety-instructed
language-model baselines across two model families, we find a large and
consistent structural-to-safe gap: the best system reaches only $0.44$
structural validity and no system exceeds $0.29$ safe-executability (95\%
CI $[0.21,0.38]$). The failure modes are complementary and quantified: the
template binds no trade parameters, the deployed generator breaks required
structure, and language models parse parameters but omit safety machinery.
On-chain, direct, constrained, and few-shot configurations each execute
$10$--$13$ swaps at roughly $33\%$ self-inflicted price impact because they
wire a slippage bound but no price-impact gate. An explicit safety
instruction is the only intervention that eliminates on-chain unsafe
execution---yet still leaves most workflows structurally incomplete---showing
the gap is real but partially addressable. We release the benchmark,
baselines, execution harness, and analysis code.
